\documentclass[12pt]{article}
\usepackage[margin=1in]{geometry}
\usepackage[T1]{fontenc}
\usepackage{amsmath}
\usepackage{bold-extra}
\usepackage{xcolor}
\definecolor{garnet}{HTML}{73000A}

\setlength{\parindent}{0.25in}
\setlength{\parskip}{0.3em}

\newcommand{\sect}[1]{\noindent{\color{garnet}\textbf{\textsc{#1}}.}}
\newcommand{\subsect}[1]{\noindent{\color{garnet}\textit{#1.}}}

\begin{document}
\sloppy

\noindent{\color{garnet}\large\textbf{\textsc{Improving Maritime Vessel Tracking and Identification Using Adaptive ST-DBSCAN Clustering}}}\\[-0.8em]
{\color{garnet}\rule{\textwidth}{0.5pt}}\\[0.2em]
\sect{Project Overview}
Density-based clustering has significantly advanced maritime radar tracking. Birant and Kut introduced ST-DBSCAN, extending DBSCAN to cluster points across both space and time, enabling coherent track association even when targets drop out intermittently \cite{birant2007}. Schubert et al. demonstrated that local density estimation substantially improves clustering in heterogeneous data, establishing theoretical foundations for adaptive parameter selection \cite{schubert2017}. In parallel, Starczewski et al. developed methods for automatic DBSCAN parameter determination using $k$-nearest neighbor statistics \cite{starczewski2020}, while Xu et al. showed that adaptive clustering is essential for handling varying sea clutter conditions in marine radar \cite{xu2021}.

However, the problem of joint spatio-temporal adaptation remains unsolved. Existing methods adapt spatial parameters or use fixed temporal windows, but maritime environments exhibit coupled spatial and temporal variability---heavy clutter causes both denser false returns and more frequent target dropouts. This project proposes a fully adaptive ST-DBSCAN framework that self-tunes both spatial radius $\varepsilon_s$ and temporal radius $\varepsilon_t$ based on local conditions, combined with vessel classification using features extracted from the resulting clusters \cite{fowdur2021}.

\sect{Research Methods}

\subsect{Problem Formalization}
ST-DBSCAN requires three parameters: spatial radius $\varepsilon_s$, temporal radius $\varepsilon_t$, and minimum points $minPts$. Standard implementations use fixed values for all parameters, but maritime environments exhibit substantial spatial and temporal variability. Heavy clutter causes more frequent target dropouts requiring larger temporal windows, while calm conditions allow tighter temporal grouping. Similarly, radar resolution varies with range, requiring spatially adaptive neighborhoods. My current approach adapts $\varepsilon_s$ using local $k$-NN distance statistics but uses fixed $\varepsilon_t = 2$ frames, which works in stable conditions but fails under varying sea states.

\subsect{Proposed Approach}
The fully adaptive ST-DBSCAN framework extends parameter estimation to the temporal dimension. Detection consistency is analyzed over recent frames to determine appropriate temporal neighborhoods: in regions where targets drop out frequently, $\varepsilon_t$ expands automatically; in stable conditions, it contracts to avoid incorrectly merging separate objects. This follows the same statistical principles used for spatial adaptation, making the entire clustering process self-tuning for both space and time.

For vessel classification, the ST-DBSCAN clusters contain rich discriminative information. Large ships produce more points, have larger spatial extent, and higher mean intensity than small boats; clutter tends to be transient while vessels persist across many frames. Features are extracted systematically across four categories: size features (point count, bounding box dimensions, convex hull area), intensity features (mean, max, and variance of radar return), temporal features (track duration, detection consistency), and motion features (velocity magnitude, heading stability). A classifier (decision tree or logistic regression) then distinguishes small boats, large ships, and clutter based on these features.

\subsect{Implementation and Evaluation}
Training data consists of labeled radar sequences collected across varying sea states using the Furuno NXT marine radar. The evaluation quantifies clustering performance using tracking accuracy metrics including track continuity and ID switches following standard MOT benchmarks \cite{bernardin2008}, temporal adaptation behavior across different clutter conditions, and classification accuracy on a held-out test set with labeled vessel types. Gap recovery using velocity-based prediction provides more continuous tracks for both clustering and classification.

\sect{Goals and Expected Results}
Goal 1: extend adaptive parameter estimation to include the temporal dimension, demonstrating self-tuning behavior across varying sea states without manual parameter adjustment. Goal 2: extract discriminative features from ST-DBSCAN clusters and train a classifier achieving at least 85\% accuracy distinguishing small boats, large ships, and clutter. Goal 3: systematic evaluation demonstrating tracking improvement across varying conditions, with quantitative comparison to fixed-parameter baselines and analysis of failure modes in heavy clutter and multi-target scenarios.

Deliverables include updated open-source code with adaptive temporal clustering and classification modules, an annotated radar dataset with vessel type labels, a technical report with quantitative evaluation, a poster presentation at the UofSC Summer Research Symposium, and a manuscript draft targeting fall 2026 conference submission.

\sect{Timeline}
The 14-week research period (350--400 hours, approximately 25--40 hours/week) proceeds as follows. Weeks 1--3: collect radar sequences across sea states, label vessel types, establish evaluation metrics. Weeks 4--7: implement adaptive temporal epsilon, test across clutter conditions, evaluate tracking improvement. Weeks 8--11: extract cluster features, analyze discriminative properties, train and test classifier. Weeks 12--14: integrate classification with tracker, Jetson benchmarking, documentation and technical report. Final week: poster preparation and symposium presentation.

\sect{NASA Relevance}
This research supports multiple NASA Mission Directorates through adaptive sensing, autonomous perception, and maritime monitoring capabilities. Under STMD, the adaptive clustering techniques address TX04 (Robotic Systems) and TX10 (Autonomous Systems) by developing self-tuning algorithms that adapt to varying environmental conditions without manual adjustment---valuable for spaceborne sensors and autonomous systems where real-time parameter tuning is impractical. The object classification component supports missions involving autonomous navigation, planetary surface analysis, and orbital debris characterization, where distinguishing object types from sensor returns is essential.

Under ARMD Strategic Thrust 6 (Assured Autonomy), robust target tracking and classification in cluttered environments establishes algorithmic foundations for autonomous systems operating in degraded conditions. The techniques developed apply to any sensor producing point clouds over time, including SAR and lidar systems relevant to NASA Earth Science missions. Additionally, NASA partners with NOAA on ocean observation; vessel classification from ground-based radar complements satellite-based ship detection and AIS tracking for comprehensive maritime domain awareness.

NASA Technology Taxonomy alignment includes TX04.4 (Mobility and manipulation for surface operations), TX10.1 (Autonomous navigation in unstructured environments), TX11.1 (Software for autonomous decision-making), and TX17.2 (Relative navigation and hazard detection).

\begin{thebibliography}{9}
\bibitem{birant2007} D. Birant and A. Kut, ``ST-DBSCAN: An Algorithm for Clustering Spatial-Temporal Data,'' \textit{Data \& Knowledge Engineering}, vol. 60, no. 1, pp. 208--221, 2007. DOI: 10.1016/j.datak.2006.01.013.
\bibitem{schubert2017} E. Schubert et al., ``DBSCAN Revisited, Revisited: Why and How You Should (Still) Use DBSCAN,'' \textit{ACM Trans. Database Syst.}, vol. 42, no. 3, pp. 1--21, 2017. DOI: 10.1145/3068335.
\bibitem{starczewski2020} A. Starczewski, P. Goetzen, and M. J. Er, ``A New Method for Automatic Determining of the DBSCAN Parameters,'' \textit{J. Artif. Intell. Soft Comput. Res.}, vol. 10, no. 3, pp. 209--221, 2020. DOI: 10.2478/jaiscr-2020-0014.
\bibitem{xu2021} Y. Xu et al., ``Adaptive Clustering-Based Marine Radar Sea Clutter Normalization,'' \textit{J. Sensors}, vol. 2021, p. 2938251, 2021. DOI: 10.1155/2021/2938251.
\bibitem{fowdur2021} J. S. Fowdur, M. Baum, and F. Heymann, ``Real-World Marine Radar Datasets for Evaluating Target Tracking Methods,'' \textit{Sensors}, vol. 21, no. 14, p. 4641, 2021. DOI: 10.3390/s21144641.
\bibitem{bernardin2008} K. Bernardin and R. Stiefelhagen, ``Evaluating Multiple Object Tracking Performance: The CLEAR MOT Metrics,'' \textit{EURASIP J. Image Video Process.}, vol. 2008, pp. 1--10, 2008. DOI: 10.1155/2008/246309.
\end{thebibliography}

\end{document}
